\section{The Bessel nonvanishing base, its parameter walk, and its literal return}
\label{sec:bessel-nonvanishing-v50}

This section isolates the original nonvanishing problem, proves its analytic
base case and its elementary parameter walk, and returns the result to the
endomorphism-valued integral actually used by the local-to-global argument.
The remaining input is identified exactly: the original two-sided
\texttt{StepStatement}, which propagates nonvanishing by one unit in a
Harish--Chandra parameter.  We do not replace that input by an abstract
branching assertion, and do not claim to have proved it here.  In particular,
the final implication to \texttt{BesselNVGap} and the supplied-form version of
\texttt{ConcreteSettingExists} is conditional on that one producer.

The proofs use DC1--DC13 for the actual large-gap discrete series and its
coefficient kernel, BWP1--BWP3 for the actual unipotent group and compact
stabilizer, and AMD8 for passage from the source's weak \texttt{IsHaarOn}
predicate to ordinary Haar measure.  They do not use a multiplicity-one
theorem, a formal-degree formula, or the desired Bessel nonvanishing itself.
Source locations and their exact pinned bytes are recorded in the accompanying
manifest.  References to a source declaration specify the target of a human
proof; the declaration is not being used in place of that proof.

\subsection{The original integral and the analytic notation}

Put $n=2r+k+1$ and use the ordered Witt decomposition
$E\oplus V\oplus F$ of BWP1, with $\dim E=\dim F=r$ and $V$ of signature
$(k,1)$.  The dual $F$ coordinates are reversed when necessary, exactly as in
the source matrix $H=H_{pk}$.  Write
\[
 G=U(H),\qquad N=N_0=\texttt{Nset}(k,r),\qquad
 H_z=\texttt{UWset}(k,r,z)\simeq U(k).
\]
Here $z\in V$ satisfies $T_z=-z^*Hz>0$.  The group $H_z$ fixes $E,F,z$
pointwise; $N$ acts trivially on the successive pieces of the prescribed
isotropic flag.  If $\ell_a$ are the original simple-root coordinates, their
last member is the $E_{r-1}$ coordinate of $nz$ and the others are the
$(a,a+1)$ entries.  For $c_a\ne0$ the original character is
\[
 \psi_{z,c}(n)=
 \exp\left(i\operatorname{Re}\sum_{a=0}^{r-1}c_a\ell_a(n)\right).
 \tag{BNV.1}
\]
For $r=0$ the sum is empty, $N=\{1\}$, and this character is $1$.

The signature $\epsilon=\epsilon^{48}$ has $k$ initial positive coordinates,
then alternating negative and positive coordinates.  Thus
$p=k+r$, $q=r+1$, and, for $k\ge1$, $q_0=\min(p,q)=r+1$.
The target parameter is
\[
 \lambda=(v_0,\ldots,v_{k-2},-b_0,-b_1,\ldots,-b_{2r+1}).
 \tag{BNV.2}
\]
A Harish--Chandra parameter of length $m$ is a strictly decreasing vector
whose coordinates belong to $(m-1)/2+\mathbb Z$.  The interlacing and
half-integrality assumptions are exactly
\[
 u_i>v_i>u_{i+1}\quad(0\le i<k-1),\qquad u_{k-1}>-b_0,
 \qquad u_i-\lambda_s\in\mathbb Z+\tfrac12.
 \tag{BNV.3}
\]

The original \texttt{BesselNVGap} has the following scope.  For every even
$k\ge2$, every $r\ge0$, every $z,c,u,v,b$ as above with $u,\lambda$
Harish--Chandra parameters and with the original \texttt{GapBig}, take
\emph{any} complete complex Hilbert space $\mathcal H$ and representation
$\sigma$ satisfying the literal predicate
$\texttt{IsDSHpk}(k,r,\epsilon,\lambda,\sigma)$.  Take any finite-dimensional
complex Hilbert space $V_\pi$ and representation $\pi$ satisfying the literal
$\texttt{IsHCTypeW}(k,r,z,u,\pi)$; take any lowest-type unit vector $w$ of
$\sigma$; and take any original \texttt{IsHaarOn} measures $\mu_W,\mu_N$ on
$H_z,N$.  The conclusion is the existence of $x,x_0\in G$ such that
\[
 \mathcal B_{\sigma,w}(x_0,x):=
 \int_{H_z\times N}
 \langle\sigma(x_0)w,\sigma(hn)\sigma(x)w\rangle
 \overline{\psi_{z,c}(n)}\,\pi(h)^{-1}\,d\mu_W(h)d\mu_N(n)
 \ne0
 \quad\text{in }\operatorname{End}_{\mathbb C}(V_\pi).
 \tag{BNV.4}
\]
Inner products are conjugate linear in the first variable.  The measures in
the formal statement live on ambient matrices and are supported on the
displayed groups; writing them on the groups makes no change.  The inverse
in (BNV.4) is the operator inverse, equal on $H_z$ to $\pi(h^{-1})$.
No Haar normalization, chosen model of $\sigma$, chosen lowest vector, or
chosen nondegenerate character may be silently substituted for these
quantifiers.

For a compact lowest type $T$ of weight
$\Lambda=\lambda+\rho_n-\rho_c$, let $\Psi_{T,\mu}$ be the normalized
coefficient kernel of DC2--DC3, with $\mu=\lambda+\rho$.  A Cayley matrix
$P$ means $P^*J_\epsilon P=H$; Euclidean unitarity of $P$ is not required.
Define the exact vector-valued source integral
\[
 A(v_1,v_2,\xi;x_0,x)=
 \int_{H_z\times N}
 \langle v_2,\Psi_{T,\mu}(Px_0^{-1}hnxP^{-1})v_1\rangle
 \overline{\psi_{z,c}(n)}\,\pi(h)^{-1}\xi\,d\mu_Wd\mu_N.
 \tag{BNV.5}
\]
This is \texttt{aaFun}, not an unspecified equivariant functional.
$\texttt{BesselNZ}(\epsilon,\lambda,u)$ means that there exist all the data
in (BNV.5), including $T,P,z,c,\pi$, the two Haar measures, and
$v_1,v_2,\xi,x_0,x$, for which (BNV.5) is nonzero.
$\texttt{BesselAll}(\epsilon,\lambda,u)$ means that for \emph{every} such
choice of data and \emph{every} $0\ne v\in T$, some $\xi,x_0,x$ make
$A(v,v,\xi;x_0,x)\ne0$.  These are the original predicates.

\subsection{A direct proof of the original subgroup-integrability input}

% BEGIN BNV1
\paragraph{BNV1. The whole original \texttt{IntSStatement}.}
\label{claim:BNV1}
For every $k,r\ge0$, admissible $z$, original Haar pair, $x_0,x\in G$, and
real $b> b_S(k,r):=2rn^2$, the function
\[
 (h,n')\longmapsto \|x_0^{-1}hn'x\|_{\rm F}^{-b}
 \tag{BNV.6}
\]
is integrable on $H_z\times N$.  Here $\|\cdot\|_{\rm F}$ is Frobenius
norm, the square root of the source's \texttt{hs}.

\begin{proof}
First suppose $r\ge1$.  In the real polynomial chart of BWP1 the dimension
is $d=2r^2+2rk+r$.  Its measure is a positive multiple of Lebesgue measure,
and the inverse-coordinate estimate of BWP2 is
\[
 1+|y(n')|\le C\|n'\|_{\rm op}^2.
\]
It follows that $\|n'\|_{\rm op}^{-b}$ is integrable whenever $b>2d$.
Indeed it is bounded by a constant times $(1+|y|)^{-b/2}$, whose radial
integral converges for $b/2>d$.  Since
\[
 2d=2r(2r+2k+1)\le2r(2r+k+1)^2=2rn^2,
\]
the asserted exponent suffices.  The elementary inequality used here holds
also when $k=0$, because $r\ge1$ and expansion gives nonnegative terms.

For $g=x_0^{-1}hn'x$,
\[
 \|n'\|_{\rm op}
 =\|h^{-1}x_0gx^{-1}\|_{\rm op}
 \le C_{x_0,x,z}\|g\|_{\rm F},
\]
uniformly for $h\in H_z$, since $H_z$ is compact.  Its Haar mass is finite.
Comparison and Tonelli give (BNV.6).  If $r=0$, $N$ is a point and the
continuous, strictly positive norm in (BNV.6) has a positive minimum on the
compact $H_z$, so the conclusion holds for every $b>0=b_S$.
AMD8 and Haar uniqueness allow every original weak Haar pair, with positive
scalars absorbed in the estimates.  Thus this is the original universal
\texttt{IntSStatement}, including its exponent and degenerate rank case.
\end{proof}
% END BNV1

\subsection{The finite parameter walk}

For this subsection let $k\ge1$, let (BNV.2) be Harish--Chandra, and assume
(BNV.3).  Write $\mathrm{Gap}_h(\eta)$ for
$|\eta_a-\eta_b|\ge h$ whenever $\epsilon_a\ne\epsilon_b$.
Suppose $b_{j+1}-b_j\ge h$.  Define
\[
 \lambda^{(t)}=
 (u_0-\tfrac12,\ldots,u_{k-1}-\tfrac12,
       -b_1-t,-b_2-2t,\ldots,-b_{2r+1}-(2r+1)t),\quad t\in\mathbb N.
 \tag{BNV.7}
\]
A predicate $\mathcal P(\eta)$ is \emph{step closed} at $(\epsilon,u,h)$ if
it is preserved by each move $\eta\mapsto\eta\pm e_s$ for which both
endpoints are Harish--Chandra and satisfy $\mathrm{Gap}_h$, and
\[
 u_j\ne\eta_s\pm\tfrac12\quad\text{for every }j.
 \tag{BNV.8}
\]
This is exactly \texttt{StepClosed}; there is no assertion here that a
particular representation-theoretic predicate has this property.

% BEGIN BNV2
\paragraph{BNV2. The sign, lattice, and interlacing arithmetic.}
\label{claim:BNV2}
The signature has $(p,q)=(k+r,r+1)$.  Every coordinate of $u$ belongs to
$(k-1)/2+\mathbb Z$, and $u$ is strictly decreasing.  Each displacement
\[
 d_i=u_i-\tfrac12-\lambda_i\quad(0\le i<k)
 \tag{BNV.9}
\]
is a nonnegative integer, where $\lambda_{k-1}=-b_0$.
Every adjacent $b$ difference is a positive integer.

\begin{proof}
Counting the alternating tail after the initial $k$ positives gives the
signature.  Interlacing gives strict decrease of $u$, and its difference
from any $\lambda_s$ belongs to $\mathbb Z+1/2$.  Since
$\lambda_s\in(n-1)/2+\mathbb Z$ and $n=k+2r+1$, this puts
$u_i$ in $n/2+\mathbb Z$, the same coset as $(k-1)/2+\mathbb Z$:
their difference is $r+1$.  Positivity and half-integrality of
$u_i-\lambda_i$ give (BNV.9).  The $-b_j$ are strictly decreasing
coordinates of the same integral coset, proving the last assertion.
\end{proof}
% END BNV2

% BEGIN BNV3
\paragraph{BNV3. Adjacent sign changes suffice for all noncompact gaps.}
\label{claim:BNV3}
If $\eta$ is strictly decreasing and
$\eta_a-\eta_{a+1}\ge h$ at every adjacent sign change, then
$\mathrm{Gap}_h(\eta)$ holds.

\begin{proof}
For $a<b$ of opposite signs, some $j$ with $a\le j<b$ changes sign.
The telescoping sum
$\eta_a-\eta_b=\sum_{i=a}^{b-1}(\eta_i-\eta_{i+1})$
is at least its $j$th positive summand, hence at least $h$.  Reversing
$a,b$ gives the absolute-value formulation.  This proves the source's full
\texttt{AdjGapStatement}, for every signature and every real $h$.
\end{proof}
% END BNV3

% BEGIN BNV4
\paragraph{BNV4. Raising the deformed tail without a forbidden crossing.}
\label{claim:BNV4}
Starting at $\lambda^{(t)}$, for $j=1,\ldots,2r+1$ in that order raise
the coordinate $-b_j-jt$ to $-b_j$ in $jt$ unit steps.  Every endpoint is
Harish--Chandra, satisfies $\mathrm{Gap}_h$, and satisfies the upward
exclusion (BNV.8).

\begin{proof}
All moves are integral, so the lattice condition is preserved.  The
preceding tail coordinate has already reached $-b_{j-1}$ when $j>1$;
the current coordinate is at most $-b_j$.  Their difference is at least
$b_j-b_{j-1}$.  The next tail coordinate is still at
$-b_{j+1}-(j+1)t$, so its difference below the current coordinate is at
least $b_{j+1}-b_j+t$.  For $j=1$ the preceding coordinate is
$u_{k-1}-1/2\ge-b_0$, giving at least $b_1-b_0$.  Thus the ordering
is strict and every adjacent sign-changing gap is at least its original
$b$ difference.  BNV3 applies.

Immediately before an upward step the moving coordinate $a$ is at most
$-b_j-1$.  In particular $a+1/2<-b_0<u_{k-1}\le u_i$, so no forbidden
equality occurs.  This proves every segment of the original
\texttt{RaiseStepStatement}, including a segment of zero length.
\end{proof}
% END BNV4

% BEGIN BNV5
\paragraph{BNV5. Lowering the initial block without a forbidden crossing.}
\label{claim:BNV5}
After BNV4, lower the initial coordinates in order
$i=k-1,k-2,\ldots,0$, taking $u_i-1/2$ to $\lambda_i$ in $d_i$ unit
steps.  All endpoints satisfy the same Harish--Chandra and gap conditions
and the downward exclusion (BNV.8).

\begin{proof}
The coordinate on the left, if any, is still $u_{i-1}-1/2$, strictly
above the current value.  Coordinates on the right have reached their
strictly smaller target values.  The lattice is unchanged.  The only
adjacent sign change involving the initial block is between positions
$k-1$ and $k$; its smallest possible difference is
$-b_0-(-b_1)=b_1-b_0$.  All other sign-changing gaps are the fixed tail
differences.  BNV3 again applies.

Before a downward step the current value $a$ obeys
$\lambda_i+1\le a\le u_i-1/2$.  If $j\le i$, then
$a-1/2\le u_i-1<u_j$.  If $j>i$, interlacing gives
$\lambda_i>u_{i+1}\ge u_j$, so
$a-1/2\ge\lambda_i+1/2>u_j$.  Thus equality is impossible in both
cases.  There are no $j>i$ when $i=k-1$.  This proves the full original
\texttt{LowerStepStatement}.
\end{proof}
% END BNV5

% BEGIN BNV6
\paragraph{BNV6. The whole original \texttt{WalkStatement}.}
\label{claim:BNV6}
For every predicate $\mathcal P$ step closed at $(\epsilon,u,h)$ and
every $t\in\mathbb N$,
\[
 \mathcal P(\lambda^{(t)})\quad\Longrightarrow\quad\mathcal P(\lambda).
 \tag{BNV.10}
\]

\begin{proof}
Apply step closure successively along the $\sum_{j=1}^{2r+1}jt$ moves of
BNV4 and then the $\sum_{i=0}^{k-1}d_i$ moves of BNV5.  These are finite
nonnegative integers.  The initial and terminal parameters are precisely
(BNV.7) and (BNV.2).  Zero-length segments need no invocation of step
closure, so the proof includes $t=0$, $r=0$, and $k=1$.
\end{proof}
% END BNV6

\subsection{The exact large-gap and large-parameter bridges}

For any signature put
\[
 \begin{split}
 h_{\rm need}(\epsilon)=\max\{&2+\max(3n/2,1+npq),\\
 &16nq_0+8\sqrt{2q_0}+2,\\
 &8\sqrt{2q_0}\,[4\max(2nq_0,2(q-1)n^2)+1]\}.
 \end{split}
 \tag{BNV.11}
\]
$\texttt{GapLap}(\epsilon,\eta,h)$ means that $\eta$ is strictly
decreasing, $\mathrm{Gap}_h(\eta)$ holds, and $h>h_{\rm need}(\epsilon)$.
The original $\texttt{StepGap}(\epsilon,r,h)$ consists of
\[
 \max(3n/2,1+npq)<h-2,\qquad
 4\max(2nq_0,2rn^2)<\frac{h}{8\sqrt{2q_0}}-1.
 \tag{BNV.12}
\]
In applications here $q_0=r+1>0$.

% BEGIN BNV7
\paragraph{BNV7. The original cubic gap implies both analytic thresholds.}
\label{claim:BNV7}
Under the original \texttt{GapBig}, there exists $h$ for which all
$b_{j+1}-b_j\ge h$, $\texttt{GapLap}(\epsilon,\lambda,h)$ holds, and
(BNV.12) holds.

\begin{proof}
The original gap supplies an odd natural $T$ such that $T\ge n+5$,
$n+T\ge1801$, and
\[
 b_{j+1}-b_j>2\bigl(n+(2r+2)T\bigr)^3+T+3.
 \tag{BNV.13}
\]
Choose $h$ to be the minimum of these finitely many differences.  Put
$Q=r+1=q=q_0$ and $B=2(n+2QT)^3$.  Since $n\ge2$ and $Q\ge1$,
$(n+5)^3\ge17n^2$; hence
\[
 B\ge16Q^3T^3\ge272Q^3n^2>160Q^2n^2.
\]
Also $\sqrt{2Q}\le2Q$ and
$\max(2nQ,2(Q-1)n^2)\le2Qn^2$.  The third member of (BNV.11) is
therefore at most
$128Q^2n^2+16Q\le132Q^2n^2$.
The second is at most $16nQ+16Q+2\le25Q^2n^2$ for $n\ge2$;
the first is at most $3+n^2Q+3n/2\le3Q^2n^2$ for $n\ge2$.
The loose bound $160Q^2n^2$ thus exceeds all three.  Equation (BNV.13)
gives $h>h_{\rm need}$, and BNV3 gives the noncompact gaps of $\lambda$.
Finally the first and third members of (BNV.11), with $q-1=r$, give
(BNV.12) directly.  The parity and $1801$ hypotheses are retained as
original input; this elementary numerical subproof needs only $T\ge n+5$.
\end{proof}
% END BNV7

% BEGIN BNV8
\paragraph{BNV8. The full large-parameter family and its lowest weights.}
\label{claim:BNV8}
Let $d_a=0$ for $a<k$ and $d_{k-1+j}=-j$ for $1\le j\le2r+1$.
For $\Lambda_t=\lambda^{(t)}+\rho_n-\rho_c$ one has
\[
 \Lambda_t=t d+\Lambda_0,\qquad
 (\Lambda_t)_a=u_a-(k-1)/2+a\quad(0\le a<k).
 \tag{BNV.14}
\]
For all sufficiently large integral $t$, $\lambda^{(t)}$ is
Harish--Chandra and
$\texttt{GapLap}(\epsilon,\lambda^{(t)},t/2)$ holds.  This is the original
\texttt{LamTStatement}, including the compact highest-weight identification.

\begin{proof}
The affine formula is immediate from (BNV.7).  At an initial positive
coordinate $a<k$, there are $r+1$ noncompact positive roots leaving $a$,
and the compact half-sum is $(r+k-1-2a)/2$.  Thus
$(\rho_n-\rho_c)_a=(2-k+2a)/2$ and (BNV.14) follows.
This is the compact $U(k)$ highest weight $u-\rho_k$.  BNV2 proves that
$u$ is itself a Harish--Chandra parameter.

The lattice condition for (BNV.7) follows from half-integrality and the
integral displacements $jt$.  The initial block is strictly decreasing.
Every remaining adjacent difference is either a fixed real number plus
$t$, or (at its boundary with the initial block) a fixed real number plus
$t$.  Consequently all such differences are positive and at least $t/2$
for all sufficiently large $t$.  These are all adjacent sign changes;
BNV3 applies.  Increasing the finite threshold so that
$t/2>h_{\rm need}(\epsilon)$ proves the assertion.  This also gives
$q_0=r+1\ge1$, as required by the original statement.
\end{proof}
% END BNV8

\subsection{An explicit Cayley matrix matched to the character}

% BEGIN BNV9
\paragraph{BNV9. The cyclic isotropic frame by Vandermonde interpolation.}
\label{claim:BNV9}
Let $x_0>\cdots>x_{2r}$ be real, put
$D_x=\operatorname{diag}(x_s)$, and equip $\mathbb R^{2r+1}$ with the form
$[a,b]=\sum_s(-1)^{s+1}a_sb_s$.  There are vectors
$e_0,\ldots,e_{r-1},z,f_0,\ldots,f_{r-1}$ with
\[
 [e_a,e_b]=[f_a,f_b]=[e_a,z]=[f_a,z]=0,\quad
 [f_a,e_b]=\delta_{ab},\quad [z,z]=-1,
\]
such that $D_xe_a=e_{a+1}$ for $a<r-1$ and $D_xe_{r-1}=z$.

\begin{proof}
Let $P_s=\prod_{j\ne s}(x_s-x_j)$, which has sign $(-1)^s$, and set
$c_s=\sqrt{(-1)^s/P_s}>0$ and $g_i=(x_s^ic_s)_s$.
Comparison of leading coefficients in the Lagrange interpolation formula
for $X^m$, $0\le m\le2r$, gives
\[
 \sum_{s=0}^{2r}\frac{x_s^m}{P_s}
 =\begin{cases}0&m<2r,\\1&m=2r.\end{cases}
\]
Hence $[g_i,g_j]=0$ if $i+j<2r$ and $[g_i,g_j]=-1$ if $i+j=2r$.
Take $e_a=g_a$ and $z=g_r$.  The Vandermonde determinant and $c_s\ne0$
show that $g_0,\ldots,g_{2r}$ form a basis.  Nondegeneracy provides
$f'_a$ with $[f'_a,g_i]=\delta_{ai}$ for $0\le i\le r$; extend these
conditions arbitrarily to the basis.  If $Q_{ab}=[f'_a,f'_b]$, replace
\[
 f_a=f'_a-\tfrac12\sum_b Q_{ab}e_b.
\]
The symmetric matrix $Q$ then gives $[f_a,f_b]=Q_{ab}-Q_{ab}/2-Q_{ba}/2=0$;
all other required pairings are unchanged.  Finally $D_xg_i=g_{i+1}$.
For $r=0$ this is simply a negative unit vector $z$, so the same
construction includes that case.
\end{proof}
% END BNV9

% BEGIN BNV10
\paragraph{BNV10. The exact trace-character identity.}
\label{claim:BNV10}
There is a Cayley matrix $P$ sending the $k$ positive middle standard
vectors to the first $k$ positive standard vectors, such that, for every
$Y\in\mathfrak n=\operatorname{Lie}(N)$,
\[
 \operatorname{tr}(D PYP^{-1})
 =\sum_{a=0}^{r-1}\bigl(\ell_a(Y)-\overline{\ell_a(Y)}\bigr),
 \quad D=\operatorname{diag}(d).
 \tag{BNV.15}
\]
With $z_0$ the standard negative middle vector and $c_a(t)=-2it$, one has
\[
 \psi_{z_0,c(t)}(\exp Y)
   =\exp\bigl(t\operatorname{tr}(D PYP^{-1})\bigr).
 \tag{BNV.16}
\]
These are the full original \texttt{PairingPropertiesStatement}, including
the source condition \texttt{WCols}.

\begin{proof}
Apply BNV9 to $x_s=-(s+1)$ in the last $2r+1$ coordinates of
$\mathbb C^n$, whose signs are $-,+,-,\ldots,-$.  Extend the frame by the
first $k$ positive standard vectors.  Use its $e_a$, these positive
vectors, $z$, and the appropriately reversed $f_a$ as columns of $P$.
Their Gram matrix is exactly $H$, so $P^*J_\epsilon P=H$, and the
positive columns have the required prescribed values.

Here is also a block verification of the trace, to fix its sign and
normalization.  Put $\Xi=P^{-1}DP$, and let $A$ send $e_a$ to
$e_{a+1}$ (or to $z$ for $a=r-1$), and vanish on the other columns.
Both $\Xi$ and $A+HA^*H$ are $H$-self-adjoint and have the same columns
on $E$ and on the positive middle space.  Their difference $R$ kills
these columns.  Self-adjointness consequently kills its rows on the
dual $F$ space and on the positive middle space.  In the block upper
nilpotent Lie algebra $\mathfrak n$, all entries that could contribute
to $\operatorname{tr}(RY)$ are therefore zero (the remaining $z,z$
entry also meets a zero diagonal entry of $Y$).

Now $\operatorname{tr}(AY)=\sum_a\ell_a(Y)$: these are just the entries
opposite the $e_a\mapsto e_{a+1},z$ arrows.  Since $HY^*H=-Y$,
\[
 \operatorname{tr}(HA^*HY)
 =-\overline{\operatorname{tr}(AY)}.
\]
This proves (BNV.15).  Each $\ell_a$ is a coordinate in the first
nilpotent filtration degree, so $(\exp Y-I)$ and $Y$ have the same
$\ell_a$ coordinates; terms $Y^j$, $j\ge2$, cannot contribute.
Finally, for $w\in\mathbb C$, $i\operatorname{Re}(-2itw)=t(w-\bar w)$.
Substituting into the original character (BNV.1) proves (BNV.16).
No unspecified scalar multiple of the character remains.
\end{proof}
% END BNV10

% BEGIN BNV11
\paragraph{BNV11. The exact exponential Haar measure on the actual $N$.}
\label{claim:BNV11}
The matrix exponential is a polynomial diffeomorphism
$\mathfrak n\longrightarrow N$.  Its inverse is the finite matrix logarithm.
The pushforward of Euclidean Lebesgue measure on $\mathfrak n$, for the
Frobenius inner product, is Haar measure on $N$, precisely the source's
$\texttt{haarLie}(N,1)$.  Moreover
\[
 \|Y\|_{\rm F}\le C\|\exp Y-I\|_{\rm F}
             (1+\|\exp Y-I\|_{\rm F})^{n-2}
 \tag{BNV.17}
\]
when $n\ge2$; if the algebra is zero the assertion is immediate.

\begin{proof}
Choose the basis order refining the isotropic flag and its reversed dual.
The endomorphisms in $\mathfrak n$ are strictly upper triangular in the
corresponding block filtration, and $N-I$ has the same property.  Thus
$Y^n=0$ and $(g-I)^n=0$.  The finite exponential and logarithm polynomials
are inverse by their formal power-series identities, valid after truncation.
The logarithm of a group element satisfies $H(\log g)^*H=-\log g$:
apply the polynomial logarithm to $Hg^*H=g^{-1}$ and use
$\log(g^{-1})=-\log g$.  The flag conditions are preserved by these
polynomials.  This proves both inclusions and the inverse assertion.
The logarithm formula
$Y=\sum_{j=1}^{n-1}(-1)^{j+1}(g-I)^j/j$
and submultiplicativity of Frobenius norm give (BNV.17), with an enlarged
constant when necessary.

In exponential coordinates multiplication is the finite
Baker--Campbell--Hausdorff polynomial.  Order coordinates by a basis
adapted to the lower central filtration
$\mathfrak n\supset[\mathfrak n,\mathfrak n]\supset\cdots$.
For fixed $Y$, the coordinate of $\log(\exp Y\exp Z)$ in filtration
degree $j$ is $Z_j+Y_j$ plus a polynomial in lower-degree coordinates
of $Z$.  Its derivative in $Z$ is block triangular with identity diagonal,
so its real Jacobian is $1$.  Change of variables proves invariance of
the pushed-forward Lebesgue measure.  It is nonzero and locally finite,
as exp is a diffeomorphism, hence Haar.  The source realification uses
exactly the real and imaginary matrix entries, whose Euclidean norm is
Frobenius norm; its prescribed pushforward with scalar $1$ is therefore
this same measure.  This proves the whole original
\texttt{NsetExpInjectiveSurjective} package, including its norm estimate.
\end{proof}
% END BNV11

% BEGIN BNV12
\paragraph{BNV12. A global Cayley majorant from unipotence and polar decomposition.}
\label{claim:BNV12}
For every Cayley matrix $P$ there is $c>0$ such that, whenever
$Y\in\mathfrak n$ and $P\exp(Y)P^{-1}=\kappa\exp X$ is its polar
decomposition in $G_\epsilon$, with $X^*=X$ and $J_\epsilon XJ_\epsilon=-X$,
\[
 c\log(1+\|Y\|_{\rm F}^2)
 \le\min\left(\|X\|_{\rm F}^2,
                  \frac{\|X\|_{\rm F}}{\sqrt{2q_0}}\right).
 \tag{BNV.18}
\]
This is the full original \texttt{CayleyMajorantBound}.

\begin{proof}
Suppose first that $q_0>0$, and write $s=\|X\|_{\rm F}$ and
$g=P\exp(Y)P^{-1}$.  Unipotence gives $\operatorname{tr}g=n$.
The nonzero eigenvalues of $X$ occur as $s_i,-s_i$, so
\[
 \begin{split}
 \|g-I\|_{\rm F}^2
 &=\operatorname{tr}(g^*g)-n
   =\operatorname{tr}(e^{2X})-n
   =4\sum_i\sinh^2s_i
 \le2s^2e^{2s}.
 \end{split}
\]
The finite logarithm of $g$ is $PYP^{-1}$.  Its polynomial formula gives
$\|Y\|_{\rm F}\le C s e^{Cs}$ for a constant depending on $P,n$.
For $0\le s\le1$, this implies
$\log(1+\|Y\|_{\rm F}^2)\le C_1s^2$; for $s\ge1$ it implies
$\log(1+\|Y\|_{\rm F}^2)\le C_2s$.
The quantity $\min(s^2,s)$ is at most
$\max(1,\sqrt{2q_0})\min(s^2,s/\sqrt{2q_0})$.
Taking a sufficiently small positive $c$ proves (BNV.18).

If $q_0=0$, $G_\epsilon$ is compact.  Its only unipotent element is $I$,
since a unitary matrix is diagonalizable.  Thus $Y=0$ and $X=0$.
The asserted inequality is $0\le0$ with the source's zero-division
convention; equivalently it can be stated without division in this
degenerate case.  No noncompact estimate is needed there.
\end{proof}
% END BNV12

\subsection{The Gaussian limit of the actual highest-weight coefficient}

For an arbitrary signature $\epsilon$ and real vectors $d,\Lambda_0$, put
$D=\operatorname{diag}(d)$ and $X^\dagger=J_\epsilon X^*J_\epsilon$.
Let $\mathcal N_+$ consist of the matrix positions $(a,b)$ such that either
$a<b$, $\epsilon_a=\epsilon_b$, $d_a\ne d_b$, or
$b<a$, $\epsilon_a\ne\epsilon_b$.
Let $\mathcal N_-$ be its transpose, and let $\mathcal L$ consist of the
positions of equal sign and equal $d$ value.  These three sets partition
all positions.  Write $Y=Y_L+Y_++Y_-$ for the corresponding projections,
and put
\[
 \tau(Y)=\operatorname{tr}(DY),\qquad
 Q(Y)=-\operatorname{tr}(D[Y,Y_-]).
 \tag{BNV.19}
\]
The subscript $-$ here refers to this explicitly specified projection;
it is not a choice of a new positive-root system.

% BEGIN BNV13
\paragraph{BNV13. Uniform highest-weight estimates from the Casimir.}
\label{claim:BNV13}
Suppose $E_{ab}$ act on an inner product space, obey the
$\mathfrak{gl}_n$ commutator relations and
$E_{ab}^*=\epsilon_a\epsilon_b E_{ba}$ on the common invariant domain.
Let $w$ be unit, compact highest of weight
$\Lambda=\lambda+\rho_n-\rho_c$, and have Casimir scalar
$\|\lambda\|^2-\|\rho\|^2$.  If $t\ge1$ and
$\Lambda=td+\Lambda_0$, then a constant $C$ depending only on
$\epsilon,d,\Lambda_0$ satisfies
\[
 \begin{aligned}
 &\|E_{ab}w\|\le C &&((a,b)\in\mathcal N_+),\\
 &\|E_{ab}w-\delta_{ab}td_aw\|\le C
                  &&((a,b)\in\mathcal L),\\
 &\left\|E(X)w-t\operatorname{tr}(DX)w\right\|
       \le C\sqrt t\sum_{a,b}|X_{ab}| &&(X\in M_n(\mathbb C)).
 \end{aligned}
 \tag{BNV.20}
\]
These are the full original \texttt{HighestWeightVectorUniformBounds}.

\begin{proof}
For $a<b$ of the same sign, compact highest weight and the commutator
give $\|E_{ba}w\|^2=\Lambda_a-\Lambda_b$.
For opposite signs the adjoint relation and the same commutator give
\[
 \|E_{ab}w\|^2=\|E_{ba}w\|^2+\Lambda_a-\Lambda_b.
\]
The Casimir identity in DC8, specialized to this highest weight, is
\[
 2\sum_{a<b,\ \epsilon_a\ne\epsilon_b}\|E_{ba}w\|^2
       =\|\rho\|^2-\|\rho_n-\rho_c\|^2=:K_0.
 \tag{BNV.21}
\]
For completeness its cancellation uses
$\Lambda+\rho_c-\rho_n=\lambda$: substitution in DC8 cancels
$\|\lambda\|^2$ from both sides and leaves exactly (BNV.21).
Thus every noncompact lower operator has a bound independent of $t$.
The compact upper operators vanish on $w$.  These give the first line
of (BNV.20).  In $\mathcal L$, compact lower operators have squared norm
$\Lambda_{0,a}-\Lambda_{0,b}$ because $d_a=d_b$; the upper ones vanish,
and diagonal operators leave residual $\Lambda_{0,a}w$.  This gives the
second line.

For every remaining off-diagonal position, the displayed norm identities
give a bound of the form
$\|E_{ab}w\|^2\le |K_0|+t|d_a-d_b|+|\Lambda_{0,a}-\Lambda_{0,b}|$.
The diagonal residuals are constant.  Since there are finitely many
positions, one constant $C$ bounds every centered operator by $C\sqrt t$.
Expand $E(X)=\sum X_{ab}E_{ab}$ and use the triangle inequality; this
proves the last line.  The argument did not assume $d$ is decreasing,
and hence has the original universal $d,\Lambda_0$ scope.
\end{proof}
% END BNV13

% BEGIN BNV14
\paragraph{BNV14. The whole original tangent-pairing estimate.}
\label{claim:BNV14}
Let $\sigma$ be any unitary representation of $G_\epsilon$ and $w$ any
smooth unit vector satisfying the highest-weight and Casimir hypotheses
of BNV13, with $\Lambda=td+\Lambda_0$ and $t\ge1$.
For every $R$ there is a constant $C_R$, depending only on
$\epsilon,d,\Lambda_0,R$, such that for every
$Y\in\mathfrak g_\epsilon$ with $\|Y\|_{\rm F}\le |R|$,
\[
 \left|
 \langle w,\sigma(\exp(Y/\sqrt t))w\rangle
       e^{-\sqrt t\tau(Y)}-e^{-Q(Y)/2}
 \right|\le\frac{C_R}{\sqrt t}.
 \tag{BNV.22}
\]
In particular the constant is uniform in the representation and its
highest vector.  This is the full original \texttt{TangentPairingStatement}.

\begin{proof}
Differentiate unitarity along real one-parameter subgroups and complexify.
The derived operators on smooth vectors then obey exactly the adjoint
relation in BNV13.  The chain rule gives, for
$U(s)=\langle w,\sigma(\exp(sY/\sqrt t))w\rangle$,
\[
 U'(s)=t^{-1/2}\langle w,\sigma(g_s)E(Y)w\rangle,
 \quad g_s=\exp(sY/\sqrt t).
\]
The conjugation rule for the derived representation and the adjoint
identity give
\[
 \langle w,\sigma(g_s)E(Y_-)w\rangle
 =\langle E((g_sY_-g_s^{-1})^\dagger)w,\sigma(g_s)w\rangle.
 \tag{BNV.23}
\]
All vectors in these identities are smooth.  Finite-dimensional matrix
Taylor expansion, with its second derivative bounded on the fixed ball,
gives uniformly for $0\le s\le1$
\[
 g_sY_-g_s^{-1}
   =Y_-+\frac{s}{\sqrt t}[Y,Y_-]+R_s,
 \qquad \|R_s\|\le C_R/t.
 \tag{BNV.24}
\]
Any fixed matrix norm can be used here; equivalence of the finite
dimensional norms only changes $C_R$.

We spell out the error estimate, since it is what permits the large-$t$
limit for a varying family of representations.  The $Y_L$ contribution
to $U'$ is $\sqrt t\tau(Y)U$ with an $O_R(t^{-1/2})$ error, by the
second line of (BNV.20); $\tau(Y_L)=\tau(Y)$.
The $Y_+$ contribution is $O_R(t^{-1/2})$ by its first line.  In
(BNV.23), $Y_-^\dagger$ belongs to $\mathcal N_+$ and gives another
$O_R(t^{-1/2})$ error.  Put $Z_s=(s/\sqrt t)[Y,Y_-]+R_s$.
Its centered action has norm at most
$C\sqrt t\|Z_s^\dagger\|=O_R(1)$, and hence contributes
$O_R(t^{-1/2})$ to $U'$.  Finally its scalar action contributes
\[
 \sqrt t\,\tau(Z_s)U(s)
   =s\operatorname{tr}(D[Y,Y_-])U(s)+O_R(t^{-1/2}).
\]
Here conjugate linearity in the first inner-product variable cancels the
conjugation in $\tau(Z_s^\dagger)=\overline{\tau(Z_s)}$.
The norm of $\sigma(g_s)w$ is $1$, so these estimates require no further
operator bounds.  Thus
\[
 U'(s)=(\sqrt t\tau(Y)-sQ(Y))U(s)+\delta_t(s),
 \qquad |\delta_t(s)|\le C_R/\sqrt t.
 \tag{BNV.25}
\]
Since $Y^\dagger=-Y$ and $D$ is real diagonal, $\tau(Y)$ is purely
imaginary.  Multiply (BNV.25) by
$\exp(-s\sqrt t\tau(Y)+s^2Q(Y)/2)$ and integrate from $0$ to $1$.
The factor involving $\tau$ has absolute value $1$, $U(0)=1$, and
$|Q(Y)|$ is bounded on the fixed ball.  This gives (BNV.22), after
enlarging $C_R$ by $\exp(\sup|Q|)$.  This derivation also proves all
differentiation, conjugation, and second-order estimates used here,
rather than treating a Gaussian approximation as an independent premise.
\end{proof}
% END BNV14

% BEGIN BNV15
\paragraph{BNV15. A nondegenerate complex Gaussian has nonzero integral.}
\label{claim:BNV15}
Let $B:E\times E\to\mathbb C$ be real bilinear on a finite-dimensional
real space and suppose $\operatorname{Re}B(y,y)>0$ for $y\ne0$.
For every additive Haar measure $dy$,
\[
 \int_E e^{-B(y,y)}\,dy\ne0.
 \tag{BNV.26}
\]

\begin{proof}
Only the symmetric part matters.  In linear coordinates write it
$A+iC$ with $A,C$ real symmetric and $A$ positive definite.
Changing variables by $A^{-1/2}$ and then an orthogonal diagonalization
of $A^{-1/2}CA^{-1/2}$ gives eigenvalues $\alpha_j\in\mathbb R$ and
\[
 \int_E e^{-B(y,y)}dy
   =c\,\pi^{d/2}(\det A)^{-1/2}
           \prod_{j=1}^{d}(1+i\alpha_j)^{-1/2},\quad c>0.
 \tag{BNV.27}
\]
The square roots are chosen continuously on the right half-plane with
positive real part.  One elementary justification of the one-dimensional
formula is as follows.  For $\operatorname{Re}z>0$, dominated
differentiation of $F(z)=\int_{\mathbb R}e^{-zx^2}dx$ and integration
by parts give $F'(z)=-F(z)/(2z)$.  Hence $F(z)\sqrt z$ is constant
on that half-plane; at $z=1$ its value is $\sqrt\pi$, by the real
Gaussian integral.  Fubini applies because the real part of the quadratic
form is positive.  Every factor in (BNV.27) is nonzero.  In dimension
zero the integral is the positive mass of the one-point Haar space.
\end{proof}
% END BNV15

% BEGIN BNV16
\paragraph{BNV16. The whole original \texttt{GaussianLimitHaar}.}
\label{claim:BNV16}
Let $E$ be a finite-dimensional real normed space with additive Haar
measure, and $f_t:E\to\mathbb C$, $t\in\mathbb N$, eventually
almost-everywhere strongly measurable.
Suppose the first assertion below is pointwise in every $y$, and the
second has one uniform threshold $t_1$: it holds for every $t\ge t_1$
and every $y\in E$:
\[
 f_t(y/\sqrt t)\longrightarrow e^{-B(y,y)},
 \qquad |f_t(y)|\le(1+c\|y\|^2)^{-\gamma t}
 \tag{BNV.28}
\]
Here $B$ is complex-valued real bilinear and $c,\gamma>0$ are fixed.
Then for all sufficiently large $t$, $f_t$ is integrable and
$\int_E f_t\ne0$.

\begin{proof}
Choose an integer $N\ge1$ beyond the eventual bounds and so large that
$2\gamma N>\dim E$.  The function $s\mapsto(1+a/s)^s$ is increasing
for $s>0$, $a\ge0$, since its logarithmic derivative is
$\log(1+a/s)-(a/s)/(1+a/s)\ge0$.
Consequently for $t\ge N$
\[
 |f_t(y/\sqrt t)|\le
       (1+(c/N)\|y\|^2)^{-\gamma N}.
 \tag{BNV.29}
\]
This is integrable by radial comparison and equivalence of norms.
Each positive scalar dilation is nonsingular for additive Haar measure,
so the scaled functions retain almost-everywhere strong measurability;
pointwise Borel measurability is not required.
Taking limits in (BNV.29) shows
\[
 e^{-\operatorname{Re}B(y,y)}
 \le(1+(c/N)\|y\|^2)^{-\gamma N}<1\quad(y\ne0).
\]
Thus the real part of the quadratic form is positive definite; no
positivity of $B$ was silently added to the original hypotheses.
Dominated convergence and BNV15 now give
\[
 \int_E f_t(y/\sqrt t)dy
   \longrightarrow \int_E e^{-B(y,y)}dy\ne0.
\]
These integrals are nonzero for all sufficiently large $t$.  Scalar
dilation of additive Haar measure gives their exact relation to the
unscaled integrals,
$\int f_t(y/\sqrt t)dy=t^{\dim E/2}\int f_t(y)dy$.
Both integrability and nonvanishing follow.  The same argument, with
dilation factor $1$, covers dimension zero.
\end{proof}
% END BNV16

\subsection{Compact orthogonality and the complete Laplace base}

% BEGIN BNV17
\paragraph{BNV17. The exact compact Schur factor in the Bessel integral.}
\label{claim:BNV17}
Let $P$ have the prescribed positive middle columns of BNV10.
Let $T$ be any compact highest-weight type whose initial $k$ weights
agree with the compact $U(k)$ highest weight of $\pi$, and let $w\in T$,
$\xi\in V_\pi$ be unit highest vectors of the corresponding weights.
Take probability Haar measure on $H_{z_0}$ and any $\sigma$-finite
measure $\mu_N$ on the ambient matrix space $M_n(\mathbb C)$, with
no support assumption.  Use the original kernel $\Psi$ on $G_\epsilon$,
extended by zero outside that closed group.  Suppose it is continuous
on $G_\epsilon$, compact bi-equivariant there, and
$m\mapsto\|\Psi(PmP^{-1})\|$ is $\mu_N$-integrable.  In this claim
$A$ means the formula (BNV.5) with its second integral over the ambient
matrix space; it agrees with the preceding notation when $\mu_N$ is
supported on $N$.  Then
\[
 \langle\xi,A(w,w,\xi;1,1)\rangle
 =\frac1{\dim V_\pi}
       \int_{M_n(\mathbb C)}\langle w,\Psi(PmP^{-1})w\rangle
                     \overline{\psi_{z_0,c}(m)}d\mu_N(m).
 \tag{BNV.30}
\]
This is the original \texttt{WhittakerSchurOrthogonality}, with its exact
normalization and its actual vector integral.

\begin{proof}
The subgroup $PH_{z_0}P^{-1}$ acts only on the first $k$ positive
coordinates.  The cyclic compact subrepresentation it generates from
$w$ is the irreducible highest-weight $U(k)$ type of $\pi$.
One can see the cyclic assertion directly from complete reducibility:
in each irreducible summand a nonzero component of $w$ is killed by the
positive-root operators and has the specified weight, so that summand
has that highest weight.  The cyclic span of this one vector is a single
diagonal copy of the type.  Finite compact highest-weight uniqueness
therefore provides an equivariant isometry $J:V_\pi\to T$ with $J\xi=w$.

For any $Y\in T$, average the rank-one operator $\xi\otimes\xi^*$
under conjugation by $\pi$.  The average commutes with the irreducible
$\pi$, hence is scalar; its trace is $1$, so the scalar is
$1/\dim V_\pi$.  Taking the appropriate matrix coefficient gives
\[
 \int_{H_{z_0}}
 \langle\xi,\pi(h)^{-1}\xi\rangle
 \langle w,T(PhP^{-1})Y\rangle,dh
      =\frac{\langle w,Y\rangle}{\dim V_\pi}.
\]
The projection of $Y$ onto $J(V_\pi)$ may be inserted without changing
either side.  Compact equivariance of the zero extension holds for
every ambient matrix: $\kappa g\in G_\epsilon$ if and only if
$g\in G_\epsilon$, so outside the group both sides are zero.
The extension is Borel measurable because the group is closed and
the kernel is continuous on it.  Moreover (BNV.1), as an exponential
of $i$ times a real number, has modulus one on every ambient matrix,
not just on $N$.  Apply the averaged identity with
$Y=\Psi(PmP^{-1})w$ for every $m$.  Unitarity bounds the norm of the
double integrand by $\|\Psi(PmP^{-1})\|$, so the hypotheses give
absolute integrability for the stated arbitrary ambient measure.
Fubini and continuity of the scalar pairing yield
(BNV.30).  Compact Haar probability exists by normalizing any Haar
measure, whose total mass is finite and positive.
\end{proof}
% END BNV17

% BEGIN BNV18
\paragraph{BNV18. The whole original \texttt{StepIBaseStatement}.}
\label{claim:BNV18}
For every $k\ge1$, $r\ge0$, and $u,v,b$ satisfying the Harish--Chandra,
interlacing, and half-integrality assumptions (BNV.2)--(BNV.3), there is
$t_0\in\mathbb N$ such that
\[
 \texttt{BesselNZ}(\epsilon,\lambda^{(t)},u)
          \quad\text{for every integer }t\ge t_0.
 \tag{BNV.31}
\]
No lower-gap assumption on the initial $b$ is added.

\begin{proof}
By BNV8, all sufficiently large $t$ have analytic gap $t/2$.
Construct the compact type $T_t$, normalized kernel $\Psi_t$, and
smooth unitary model of DC1--DC11 for this parameter.  After dividing
the embedding by the square root of the positive scalar in DC11,
the unit highest vector $w_t\in T_t$ gives the exact coefficient
\[
 \langle w_t,\Psi_t(g)w_t\rangle
    =\langle\widetilde w_t,\sigma_t(g)\widetilde w_t\rangle,
\]
where $\widetilde w_t$ is smooth, unit, of lowest compact highest weight
$\Lambda_t$, and has the required Casimir.  The gap hypothesis here
is stronger than DC's construction threshold by (BNV.11).

Fix the one matrix $P$ of BNV10, independently of $t$.
On the real vector space $\mathfrak n$ define
\[
 f_t(Y)=\langle w_t,\Psi_t(P\exp YP^{-1})w_t\rangle
       \exp\bigl(-t\operatorname{tr}(DPYP^{-1})\bigr).
 \tag{BNV.32}
\]
The phase has modulus $1$.  By BNV10 it is precisely the conjugate
of the character $\psi_{z_0,c(t)}(\exp Y)$ with every
$c_a(t)=-2it\ne0$ for $t\ge1$.  Applying BNV14 to $PYP^{-1}$ gives
\[
 f_t(Y/\sqrt t)\longrightarrow
     \exp\bigl(-Q(PYP^{-1})/2\bigr).
 \tag{BNV.33}
\]
The exponent is the diagonal of a real bilinear form on $\mathfrak n$:
one may take
$B(Y,Z)=-\operatorname{tr}(D[PYP^{-1},(PZP^{-1})_-])/2$.
All its data are fixed as $t$ varies.

For $P\exp YP^{-1}=\kappa e^X$, DC6 and BNV12 give the global estimate
\[
 \begin{split}
 \|\Psi_t(P\exp YP^{-1})\|
 &\le\exp\left[-\frac{t}{16}
       \min(\|X\|_{\rm F}^2,\|X\|_{\rm F}/\sqrt{2q_0})\right]\\
 &\le(1+\|Y\|_{\rm F}^2)^{-ct/16}.
 \end{split}
 \tag{BNV.34}
\]
This bound is uniform in $t$; in particular no dimension of $T_t$
occurs in it.  BNV16 applied to (BNV.32)--(BNV.34) proves that the
scalar integral of $f_t$ is nonzero for every sufficiently large $t$.
The same bound makes the full operator norm in (BNV.34) integrable
once $ct/8>\dim\mathfrak n$.

Choose $\mu_N$ to be the exact exponential Haar measure of BNV11,
$\mu_W$ the probability Haar on $H_{z_0}$, and $\pi$ the compact
$U(k)$ type with Harish--Chandra parameter $u$, realized on $H_{z_0}$.
BNV2 gives its integral highest weight, and BNV8 identifies that
weight with the initial $k$ entries of $\Lambda_t$.
Choose a unit highest vector $\xi$.  BNV17 gives
\[
 \langle\xi,A(w_t,w_t,\xi;1,1)\rangle
       =\frac1{\dim V_\pi}\int_{\mathfrak n}f_t(Y)dY\ne0.
\]
All these are legitimate witnesses for the original \texttt{BesselNZ}:
the character is nondegenerate, the Cayley matrix has the prescribed
form, the measures have the original Haar predicates, and both group
points are $1$.  This proves (BNV.31) with every quantifier of the
source base statement.  If $r=0$, $\mathfrak n$ is zero-dimensional;
the same argument reduces to the nonzero compact Schur factor, with
no character coordinates to check.
\end{proof}
% END BNV18

\subsection{From one nonzero model period to every original datum}

The next four claims are proved for a Harish--Chandra parameter $\eta$,
a signature with $(p,q)=(k+r,r+1)$, $k\ge1$, and some
$\texttt{GapLap}(\epsilon,\eta,h)$.  They do not require that $\eta$
have the special form (BNV.2).  In this scope $q_0=r+1$.

% BEGIN BNV19
\paragraph{BNV19. Absolute convergence for every original coefficient period.}
\label{claim:BNV19}
For all data permitted in (BNV.5), every $v_1,v_2,\xi$ and every
$x_0,x\in G$, its integrand is Bochner integrable.  The analogous
endomorphism integrand is also integrable.  This holds for every original
Haar pair and every Cayley matrix $P$.

\begin{proof}
Let $c_0=\sqrt{2q_0}$ and $a=h/(8c_0)$.  In polar coordinates
$g=\kappa e^X$, $s=\|X\|_{\rm F}$, DC6 gives
\[
 \|\Psi(g)\|\le e^{-(h/8)\min(s^2,s/c_0)}
                 \le C e^{-as}.
\]
Indeed $s^2\ge s/c_0-1/(4c_0^2)$, so one may take
$C=e^{h/(32c_0^2)}$.  Also $\|g\|_{\rm F}\le\sqrt n e^s$.
Consequently
\[
 \|\Psi(PgP^{-1})\|\le C_P\|g\|_{\rm F}^{-a}.
 \tag{BNV.35}
\]
Conjugation by a fixed invertible $P$ changes Frobenius norm only by
fixed positive comparison constants.  The third member of (BNV.11)
implies
\[
 a>4\max(2nq_0,2rn^2)+1>2rn^2=b_S.
\]
Apply BNV1 to $g=x_0^{-1}hn'x$.  On the measure's support,
$|\psi(n')|=1$ and $\|\pi(h)^{-1}\xi\|=\|\xi\|$, while
$\|\pi(h)^{-1}\|=1$.  Thus each asserted norm is bounded by a fixed
constant times the integrable function in BNV1.  Continuity on the
groups gives strong measurability; finite-dimensional ranges of the
kernel and of $\pi$ suffice for the operator-valued assertion.
This proves actual absolute convergence, including every translation
and every subsequent finite linear combination used below.
\end{proof}
% END BNV19

% BEGIN BNV20
\paragraph{BNV20. Explicit conjugacy of all allowed Bessel characters and compact types.}
\label{claim:BNV20}
For any two allowed pairs $(z_1,c_1)$ and $(z_2,c_2)$ there is
$a\in G$ normalizing $N$ such that
\[
 aH_{z_1}a^{-1}=H_{z_2},\qquad
 \psi_{z_2,c_2}(an'a^{-1})=\psi_{z_1,c_1}(n')\quad(n'\in N).
 \tag{BNV.36}
\]
The transport preserves the compact Harish--Chandra label $u$.
Thus if $\pi_1,\pi_2$ have that label, there is an isometry
$J:V_{\pi_1}\to V_{\pi_2}$ intertwining
$\pi_1(h)$ with $\pi_2(aha^{-1})$.

\begin{proof}
Write $T_i=-z_i^*Hz_i$ and $s=\sqrt{T_1/T_2}>0$.
The negative unit vectors $z_i/\sqrt{T_i}$ have positive-definite
orthogonal complements in $V$.  Choosing orthonormal bases of those
complements constructs $M\in U(k,1)$ with $Mz_1=sz_2$.
Let $D_E=\operatorname{diag}(d_0,\ldots,d_{r-1})$, and use the Levi
element
\[
 a=\operatorname{diag}(D_E,M,D_E^{-*}),
\]
with the dual ordering adjusted as for $H$.  It normalizes the isotropic
flag and $N$.  For $r>0$ choose
\[
 d_{r-1}=s\,c_{1,r-1}/c_{2,r-1},\qquad
 d_j=(c_{1,j}/c_{2,j})d_{j+1}\quad(j=r-2,\ldots,0).
\]
All are nonzero.  Direct entry multiplication gives
\[
 \ell_j(an'a^{-1})=(d_j/d_{j+1})\ell_j(n')\quad(j<r-1),
 \qquad
 \ell_{r-1,z_2}(an'a^{-1})=(d_{r-1}/s)\ell_{r-1,z_1}(n').
\]
Thus $c_{2,j}\ell_{j,z_2}(an'a^{-1})=c_{1,j}\ell_{j,z_1}(n')$
for every $j$, proving the character identity exactly, before taking
real parts or exponentials.  The fixed spaces of $H_{z_i}$ also
transform as asserted.  For $r=0$ only $M$ is needed.

For the label assertion choose a positive orthonormal frame in
$z_1^\perp\cap V$; its image under $M$ is such a frame for $z_2$.
In these frames conjugation by $a$ is the identity identification of
the compact $U(k)$ factors, because $D_E$ acts on the fixed outer
coordinates.  Another allowed frame differs by an element of $U(k)$,
which induces an inner automorphism and does not change the compact
type.  Finite compact highest-weight uniqueness supplies the unitary
intertwiner $J$.  This checks the frame clause of the original
$\texttt{IsHCTypeW}$, rather than presuming that an arbitrary group
isomorphism preserves its label.
\end{proof}
% END BNV20

% BEGIN BNV21
\paragraph{BNV21. The full original transport $\texttt{BesselNZ}\Rightarrow\texttt{BesselAll}$.}
\label{claim:BNV21}
Under the analytic gap specified above,
\[
 \texttt{BesselNZ}(\epsilon,\eta,u)
       \Longrightarrow\texttt{BesselAll}(\epsilon,\eta,u).
 \tag{BNV.37}
\]
All choices in the latter predicate, including every nonzero vector of
the lowest compact type, are retained.

\begin{proof}
Start with any nonzero witness (BNV.5).  We make each change of data
explicit.  All integrals in the argument converge by BNV19.

First, compact highest-weight uniqueness identifies any two choices
of $T$ unitarily, carrying their highest lines to each other.  The
normalized compact average defining $\Psi$ in DC2 is then conjugated
by that isometry: its rank-one highest-line projector and its
representation are conjugated, and its dimension and Haar normalization
are unchanged.  Also $\mu$ is forced to be $\eta+\rho$.  Thus neither
choice changes the existence of a nonzero period.

Second, for two Cayley matrices $P,P'$ write $P'=Pq$; the identity
$P^*J_\epsilon P=P'^*J_\epsilon P'=H$ gives $q\in G$.
Direct substitution in (BNV.5) gives
\[
 A_{P'}(v_1,v_2,\xi;x_0q,xq)
           =A_P(v_1,v_2,\xi;x_0,x).
 \tag{BNV.38}
\]
No invariance of the subgroup under $q$ is being claimed or needed;
only the two free group points have changed.

Third, use $a$ from BNV20 to change $z,c$, and its isometry $J$ to
change the compact representation.  By AMD8 and uniqueness of Haar
measure, the target Haar measures are positive multiples
$\alpha,\beta$ of the pushforwards of the old measures under
$h\mapsto aha^{-1}$ and $n'\mapsto an'a^{-1}$.
Changing variables gives the exact vector identity
\[
 A_2(v_1,v_2,J\xi;ax_0,ax)
       =\alpha\beta\,J A_1(v_1,v_2,\xi;x_0,x).
 \tag{BNV.39}
\]
The matrix argument reduces exactly:
\[
 (ax_0)^{-1}(aha^{-1})(an'a^{-1})(ax)=x_0^{-1}hn'x.
\]
The character is (BNV.36), and the representation factor satisfies
$\pi_2(aha^{-1})^{-1}J=J\pi_1(h)^{-1}$.
This also allows arbitrary Haar normalizations.  We have now obtained
a nonzero period for the entire target datum, still with possibly
different vectors $v_1,v_2$.

Finally fix any nonzero $v\in T$.  The finite-dimensional span of its
$K_\epsilon$ orbit is a nonzero invariant subspace, hence all of the
irreducible compact type $T$.  Express $v_1$ and $v_2$ as finite linear
combinations of $T(k_1)v$ and $T(k_2)v$.  Compact bi-equivariance gives
\[
 \begin{split}
 &\langle T(k_2)v,\Psi(Px_0^{-1}hn'xP^{-1})T(k_1)v\rangle\\
 &\qquad=\langle v,
 \Psi\bigl(P(x_0P^{-1}k_2P)^{-1}hn'(xP^{-1}k_1P)P^{-1}\bigr)v\rangle.
 \end{split}
\]
Consequently the known nonzero vector period is a finite linear
combination of periods with the same vector $v$ in both slots and
with new free group points.  At least one term is nonzero, proving
(BNV.37).  This proves the original universal transport by explicit
integrable changes of data; it introduces no derivative-decay or
uniqueness-of-Bessel-model premise.
\end{proof}
% END BNV21

% BEGIN BNV22
\paragraph{BNV22. Return to every original representation and its endomorphism integral.}
\label{claim:BNV22}
Suppose $\texttt{BesselAll}(\epsilon,\eta,u)$ holds in the analytic-gap
scope.  Then for every original $\sigma,w,\pi,z,c,\mu_W,\mu_N$ permitted
in (BNV.4), some $x_0,x\in G$ make the \emph{actual endomorphism-valued}
integral (BNV.4) nonzero.

\begin{proof}
First reconcile the independently existential unitary Cayley witnesses
in the original \texttt{IsDSHpk} and \texttt{IsLowUnitH} predicates.
If they are $P_0,P_1$, then $P_1P_0^{-1}$ is both Euclidean unitary
and $J_\epsilon$-unitary, hence belongs to $K_\epsilon$.
The resulting inner compact conjugacy preserves the entire lowest
$K_\epsilon$-isotypic subspace.  Thus the original $w$ belongs to the
lowest type for the same Cayley witness used for the representation.
The isomorphism theorem DC13 identifies the arbitrary supplied
$\texttt{IsDSHpk}$ representation with the constructed large-gap model,
after its allowed Cayley change.  It identifies its entire lowest
compact type with $T$.  The coefficient identity DC11, with the
embedding normalized isometrically, therefore gives a nonzero $v\in T$
such that
\[
 \langle w,\sigma(g)w\rangle
       =\langle v,\Psi(PgP^{-1})v\rangle\qquad(g\in G).
 \tag{BNV.40}
\]
Equivalently one may keep DC11's positive scalar; it multiplies all
following integrals by the same nonzero number.  No choice of the
original unit vector $w$ is being imposed.

By $\texttt{BesselAll}$ choose $\xi,x_0,x$ with
$A(v,v,\xi;x_0,x)\ne0$ for the original compact type, character, and
Haar pair.  Unitarity gives
\[
 \langle\sigma(x_0)w,\sigma(hn')\sigma(x)w\rangle
 =\langle w,\sigma(x_0^{-1}hn'x)w\rangle.
\]
Substitute (BNV.40).  BNV19 proves integrability of the operator
integrand, and evaluation $L\mapsto L\xi$ is a continuous linear map
on its finite-dimensional endomorphism space.  It commutes with the
Bochner integral, giving
\[
 \mathcal B_{\sigma,w}(x_0,x)\xi=A(v,v,\xi;x_0,x)\ne0.
\]
Thus the operator itself is nonzero.  This is precisely the original
conclusion, including its order of factors $hn'$, conjugated character,
two group translates, arbitrary lowest unit vector, and every original
Haar measure.  A nonzero abstract functional has not been substituted
for that conclusion.
\end{proof}
% END BNV22

\subsection{The remaining producer and the conditional original endpoint}

Here is the remaining source input, with its full scope.  The original
\texttt{StepStatement} says that, assuming the original
\texttt{IntSStatement}, for every $k,r$, every signature $\epsilon$ with
$p=k+r$, $q=r+1$, every $u$ and every real $h$ satisfying (BNV.12),
the predicate
\[
 \eta\longmapsto\texttt{BesselNZ}(\epsilon,\eta,u)
 \tag{BNV.41}
\]
is step closed in exactly the sense (BNV.8).  In expanded form it has
two assertions: if $\eta$ and $\eta-e_s$ are Harish--Chandra and have
noncompact gaps at least $h$, if $u_j\ne\eta_s-1/2$ for every $j$, and
if the period for $\eta$ is nonzero, then the period for $\eta-e_s$ is
nonzero; and the analogous assertion for $\eta+e_s$ with exclusion
$u_j\ne\eta_s+1/2$.  The data witnessing the new existential period may
change.  This statement is a retained, unreconstructed source producer
in this section, not an external published theorem or a newly claimed
proof obligation closure.

The source proves the two directions in \texttt{LaplaceEstimates/StepDown}
and \texttt{StepUp}.  Their substantive inputs concern standard and dual
tensor products, complementary Capelli minors, quotients with the shifted
Casimir and lowest weight, and nonzero period evaluation.  The factors
$\prod_j(u_j-\eta_s\pm1/2)$ explain the exclusions, but merely writing
these factors does not prove that the required quotient carries a
nonzero period.  That tensor-and-quotient argument is the remaining
analytic reconstruction task.  BNV1 removes its original subgroup
integrability antecedent; BNV18 removes its large-parameter base
antecedent; BNV6 supplies the entire finite path that will consume it.

% BEGIN BNV23
\paragraph{BNV23. Exact conditional closure of \texttt{BesselNVGap} and the supplied-form setting.}
\label{claim:BNV23}
The original \texttt{StepStatement}, together with the inherited human
inputs named above, implies the whole original \texttt{BesselNVGap}.
On the same supplied $\texttt{Forms0}$ scope as ASG, AMD, and BWP24, it
therefore implies the original \texttt{ConcreteSettingExists}, with
that supplied-form premise retained.

\begin{proof}
Take arbitrary data in the original target (BNV.4).  BNV7 supplies
$h$ satisfying the original tail gaps, $\texttt{GapLap}$, and
$\texttt{StepGap}$.  BNV1 proves \texttt{IntSStatement}, so the retained
\texttt{StepStatement} makes (BNV.41) step closed.  By BNV18 choose
an integer $t\ge t_0$ for which
$\texttt{BesselNZ}(\epsilon,\lambda^{(t)},u)$ holds.
BNV6 then gives $\texttt{BesselNZ}(\epsilon,\lambda,u)$.
BNV21 gives $\texttt{BesselAll}(\epsilon,\lambda,u)$, and BNV22 gives
the nonzero endomorphism integral for the original arbitrary data.
This proves every quantifier of \texttt{BesselNVGap} under the one
explicitly retained producer.

For the global setting, invoke the already proved BWP24.  Its input is
this very \texttt{BesselNVGap}, not a branching substitute.  BWP24
constructs the local wave packet from the original coefficient and
nonvanishing integrals, transports it with the original frame and
$4\pi/g$ character, normalizes the global block before any later
\texttt{MBlk}, and joins ASG and AMD on the same data and exact measures.
It yields \texttt{ConcreteSettingExists} for each supplied block sequence
in their original \texttt{Forms0} scope.  No new choice of ambient forms,
no strengthening of that scope, and no orbital/fiber or Siegel--Weil
comparison follows from this invocation.
\end{proof}
% END BNV23

\paragraph{Resulting boundary.}
This section completely reconstructs the original subgroup integrability,
parameter walk, large-parameter Laplace base, universal-data transport,
and the return from a coefficient period to the prescribed endomorphism
integral.  The two one-unit propagation directions remain to be
reconstructed.  Thus (BNV.4), the supplied-form global setting, and all
subsequent orbital/fiber comparisons remain conditional at this stage.
No whole selected obligation, unrestricted RH theorem, or unrestricted
Dirichlet-GRH theorem is declared closed by this component.
